\textbf{Phần 1. Trắc nghiệm lựa chọn}
\setcounterpageref{ex}{0}

\begin{ex} %[3H9N52]
Muốn tính diện tích hình chữ nhật ta lấy
\choice
{(Chiều dài + Chiều rộng) $\times 2$}
{\True Chiều dài nhân với chiều rộng (cùng đơn vị đo)}
{Chiều dài nhân với 4}
{Chiều dài cộng chiều rộng}
\loigiai{
Diện tích hình chữ nhật = Chiều dài $\times$ Chiều rộng.\\
\textbf{Đáp án của câu hỏi là B}
}
\end{ex}

\begin{ex} %[3H9N52]
Muốn tính diện tích hình vuông ta lấy
\choice
{\True Độ dài một cạnh nhân với chính nó}
{Độ dài cạnh nhân với 4}
{Độ dài cạnh nhân với 2}
{Độ dài cạnh cộng với chính nó}
\loigiai{
Diện tích hình vuông = Cạnh $\times$ Cạnh.\\
\textbf{Đáp án của câu hỏi là A}
}
\end{ex}

\begin{ex} %[3H9H52]
Một hình chữ nhật có chiều dài $8\text{ cm}$, chiều rộng $5\text{ cm}$. Diện tích hình chữ nhật là
\choice
{$26\text{ cm}^2$}
{\True $40\text{ cm}^2$}
{$35\text{ cm}^2$}
{$45\text{ cm}^2$}
\loigiai{
$8 \times 5 = 40\text{ cm}^2$.\\
\textbf{Đáp án của câu hỏi là B}
}
\end{ex}

\begin{ex} %[3H9H52]
Một hình vuông có cạnh dài $6\text{ cm}$. Diện tích hình vuông là
\choice
{$24\text{ cm}^2$}
{$30\text{ cm}^2$}
{$12\text{ cm}^2$}
{\True $36\text{ cm}^2$}
\loigiai{
$6 \times 6 = 36\text{ cm}^2$.\\
\textbf{Đáp án của câu hỏi là D}
}
\end{ex}

\begin{ex} %[3H9H52]
Một hình vuông có cạnh dài $8\text{ cm}$. Diện tích hình vuông là
\choice
{$32\text{ cm}^2$}
{$48\text{ cm}^2$}
{\True $64\text{ cm}^2$}
{$56\text{ cm}^2$}
\loigiai{
$8 \times 8 = 64\text{ cm}^2$.\\
\textbf{Đáp án của câu hỏi là C}
}
\end{ex}

\begin{ex} %[3H9H52]
Một hình chữ nhật có chiều dài $12\text{ cm}$, chiều rộng $6\text{ cm}$. Diện tích hình chữ nhật là
\choice
{\True $72\text{ cm}^2$}
{$36\text{ cm}^2$}
{$60\text{ cm}^2$}
{$84\text{ cm}^2$}
\loigiai{
$12 \times 6 = 72\text{ cm}^2$.\\
\textbf{Đáp án của câu hỏi là A}
}
\end{ex}

\begin{ex} %[3H9V52]
Một hình vuông có chu vi $28\text{ cm}$. Diện tích hình vuông đó là
\choice
{$28\text{ cm}^2$}
{$56\text{ cm}^2$}
{\True $49\text{ cm}^2$}
{$64\text{ cm}^2$}
\loigiai{
Cạnh hình vuông: $28 : 4 = 7\text{ cm}$.\\
Diện tích: $7 \times 7 = 49\text{ cm}^2$.\\
\textbf{Đáp án của câu hỏi là C}
}
\end{ex}

\begin{ex} %[3H9V52]
Một hình chữ nhật có diện tích $45\text{ cm}^2$, chiều dài là $9\text{ cm}$. Chiều rộng của hình chữ nhật là
\choice
{$4\text{ cm}$}
{$6\text{ cm}$}
{$7\text{ cm}$}
{\True $5\text{ cm}$}
\loigiai{
Chiều rộng: $45 : 9 = 5\text{ cm}$.\\
\textbf{Đáp án của câu hỏi là D}
}
\end{ex}

\begin{ex} %[3H9V52]
Một hình chữ nhật có chiều rộng $4\text{ cm}$, chiều dài gấp 3 lần chiều rộng. Diện tích hình chữ nhật đó là
\choice
{\True $48\text{ cm}^2$}
{$36\text{ cm}^2$}
{$52\text{ cm}^2$}
{$60\text{ cm}^2$}
\loigiai{
Chiều dài: $4 \times 3 = 12\text{ cm}$.\\
Diện tích: $12 \times 4 = 48\text{ cm}^2$.\\
\textbf{Đáp án của câu hỏi là A}
}
\end{ex}

\begin{ex} %[3H9V52]
Một tấm thảm hình vuông có cạnh $9\text{ dm}$. Diện tích tấm thảm là
\choice
{$36\text{ dm}^2$}
{$72\text{ dm}^2$}
{$90\text{ dm}^2$}
{\True $81\text{ dm}^2$}
\loigiai{
$9 \times 9 = 81\text{ dm}^2$.\\
\textbf{Đáp án của câu hỏi là D}
}
\end{ex}

\begin{ex} %[3H9V52]
Một tấm bìa hình chữ nhật có chiều dài $15\text{ cm}$, chiều rộng $8\text{ cm}$. Người ta cắt đi một hình vuông cạnh $5\text{ cm}$. Diện tích phần bìa còn lại là
\choice
{$90\text{ cm}^2$}
{$100\text{ cm}^2$}
{\True $95\text{ cm}^2$}
{$85\text{ cm}^2$}
\loigiai{
Diện tích chữ nhật: $15 \times 8 = 120\text{ cm}^2$.\\
Diện tích cắt đi: $5 \times 5 = 25\text{ cm}^2$.\\
Còn lại: $120 - 25 = 95\text{ cm}^2$.\\
\textbf{Đáp án của câu hỏi là C}
}
\end{ex}

\begin{ex} %[3H9C52]
Một hình vuông có diện tích $36\text{ cm}^2$. Chu vi của hình vuông đó là
\choice
{$18\text{ cm}$}
{\True $24\text{ cm}$}
{$30\text{ cm}$}
{$36\text{ cm}$}
\loigiai{
Vì $6 \times 6 = 36$ nên cạnh hình vuông là $6\text{ cm}$.\\
Chu vi: $6 \times 4 = 24\text{ cm}$.\\
\textbf{Đáp án của câu hỏi là B}
}
\end{ex}

\textbf{Phần 2. Trắc nghiệm đúng sai}
\setcounterpageref{ex}{0}

\begin{ex} %[3H9H52]
Xét tính đúng sai về diện tích:
\choiceTF
{\True Diện tích hình vuông cạnh $5\text{ cm}$ là $25\text{ cm}^2$.}
{\True Diện tích hình chữ nhật dài $6\text{ cm}$, rộng $4\text{ cm}$ là $24\text{ cm}^2$.}
{Hình vuông có chu vi $16\text{ cm}$ thì có diện tích $16\text{ cm}$.}
{\True Hình vuông có chu vi $16\text{ cm}$ thì có diện tích $16\text{ cm}^2$.}
\loigiai{
\begin{itemize}
    \item $5 \times 5 = 25\text{ cm}^2$: Đúng.
    \item $6 \times 4 = 24\text{ cm}^2$: Đúng.
    \item Đơn vị diện tích phải là $\text{cm}^2$ (không phải $\text{cm}$). Ý 3 Sai, Ý 4 Đúng.
\end{itemize}
\begin{itemize}
    \item a) \textbf{Đúng}.
    \item b) \textbf{Đúng}.
    \item c) \textbf{Sai}.
    \item d) \textbf{Đúng}.
\end{itemize}
}
\end{ex}

\begin{ex} %[3H9V52]
Cho hình chữ nhật có chu vi $28\text{ cm}$, chiều dài là $8\text{ cm}$. Xét tính đúng sai:
\choiceTF
{\True Nửa chu vi hình chữ nhật là $14\text{ cm}$.}
{\True Chiều rộng hình chữ nhật là $6\text{ cm}$.}
{Diện tích hình chữ nhật là $42\text{ cm}^2$.}
{\True Diện tích hình chữ nhật là $48\text{ cm}^2$ ($8 \times 6 = 48$).}
\loigiai{
\begin{itemize}
    \item Nửa chu vi: $28 : 2 = 14\text{ cm}$. Đúng.
    \item Rộng: $14 - 8 = 6\text{ cm}$. Đúng.
    \item Diện tích: $8 \times 6 = 48\text{ cm}^2 \neq 42\text{ cm}^2$. Ý 3 Sai, Ý 4 Đúng.
\end{itemize}
\begin{itemize}
    \item a) \textbf{Đúng}.
    \item b) \textbf{Đúng}.
    \item c) \textbf{Sai}.
    \item d) \textbf{Đúng}.
\end{itemize}
}
\end{ex}

\textbf{Phần 3. Trả lời ngắn}
\setcounterpageref{ex}{0}

\begin{ex} %[3H9H52]
Tính diện tích hình chữ nhật có chiều dài $9\text{ cm}$, chiều rộng $7\text{ cm}$ theo xăng-ti-mét vuông.
\shortans{63}
\loigiai{
\textbf{Đáp số:} $63$\\
$9 \times 7 = 63\text{ cm}^2$.
}
\end{ex}

\begin{ex} %[3H9H52]
Tính diện tích hình vuông có cạnh dài $8\text{ cm}$ theo xăng-ti-mét vuông.
\shortans{64}
\loigiai{
\textbf{Đáp số:} $64$\\
$8 \times 8 = 64\text{ cm}^2$.
}
\end{ex}

\begin{ex} %[3H9V52]
Một hình vuông có chu vi $36\text{ cm}$. Tính diện tích hình vuông đó theo xăng-ti-mét vuông.
\shortans{81}
\loigiai{
\textbf{Đáp số:} $81$\\
Cạnh: $36 : 4 = 9\text{ cm}$. Diện tích: $9 \times 9 = 81\text{ cm}^2$.
}
\end{ex}

\begin{ex} %[3H9V52]
Một hình chữ nhật có diện tích $54\text{ cm}^2$, chiều rộng $6\text{ cm}$. Tính chiều dài hình chữ nhật theo xăng-ti-mét.
\shortans{9}
\loigiai{
\textbf{Đáp số:} $9$\\
$54 : 6 = 9\text{ cm}$.
}
\end{ex}

\begin{ex} %[3H9V52]
Một hình chữ nhật có chiều rộng $5\text{ cm}$, chiều dài $12\text{ cm}$. Tính diện tích hình chữ nhật theo xăng-ti-mét vuông.
\shortans{60}
\loigiai{
\textbf{Đáp số:} $60$\\
$12 \times 5 = 60\text{ cm}^2$.
}
\end{ex}

\begin{ex} %[3H9C52]
Một hình chữ nhật có chu vi $20\text{ cm}$ và có chiều dài gấp 4 lần chiều rộng. Tính diện tích hình chữ nhật theo xăng-ti-mét vuông.
\shortans{16}
\loigiai{
\textbf{Đáp số:} $16$\\
Nửa chu vi: $20 : 2 = 10\text{ cm}$.\\
Rộng: $10 : 5 = 2\text{ cm}$, Dài: $8\text{ cm}$.\\
Diện tích: $8 \times 2 = 16\text{ cm}^2$.
}
\end{ex}
